
Fine-tuning a large language model requires allocating parameters to each layer. In practice this allocation is made uniformly, before training begins, despite the fact that individual layers contribute differently to model behavior. Low-Rank Adaptation (LoRA; Hu et al., 2022) constrains weight updates to low-rank matrices, reducing trainable parameters by orders of magnitude. The rank r of these updates is set uniformly across all transformer layers --- a choice that does not account for the heterogeneous computational roles that emerge during pre-training.

Aghajanyan et al. (2021) demonstrate that pre-trained language models have low intrinsic dimension for fine-tuning: more than 90\% of full fine-tuning performance can be recovered in a subspace far smaller than the full parameter space. Structure-aware projections (SAID) significantly outperform dimension-insensitive alternatives, suggesting that layers differ in their fine-tuning dimensionality. If this is the case, allocating equal rank to all layers may waste parameters in low-complexity layers while providing insufficient capacity in high-complexity ones.

Adaptive rank methods have addressed this concern. AdaLoRA \citep{zhang2023adalora} applies singular value decomposition to weight update matrices and prunes singular values during training via an importance score. DyLoRA \citep{valipour2022dylora} trains across multiple rank values simultaneously through stochastic rank sampling. Both methods share a fundamental constraint: the rank allocation is learned during training, which introduces a circular dependency --- the allocation is determined by the training process it is intended to guide --- and requires a full training run per task.

This work investigates whether the pre-trained model itself carries a signal that could identify which layers require more rank, before any fine-tuning begins. Specifically, the layer-wise MLP activation sparsity profile of LLaMA-3.1-8B is measured and characterized across three dimensions:

\begin{enumerate}
\item \textbf{Heterogeneity (RQ1):} Do layers exhibit significantly different activation sparsity (CV > 0.3) with a stable cross-distribution ranking?
\item \textbf{Cross-distribution stability (RQ2):} Is the sparsity profile stable across diverse calibration text distributions?
\item \textbf{Threshold invariance (RQ3):} Is the layer rank ordering invariant to the choice of epsilon threshold across two orders of magnitude?
\end{enumerate}

The central finding is that LLaMA-3.1-8B's sparsity profile passes all three tests by substantial margins, suggesting that this profile reflects pre-training geometry rather than input-dependent activation patterns. The practical implication --- whether this profile can guide rank allocation and achieve efficiency gains --- is designated future work (H-M3, H-M4) and was not executed in this study.

\textbf{Contributions:}

\begin{enumerate}
\item The first systematic characterization of LLaMA-3.1-8B's layer-wise MLP activation sparsity profile: CV = 0.544 (epsilon = 0.01) with a systematic depth gradient (early layers highest sparsity, deep layers lowest).
\item Cross-distribution stability: ICC(3,k) = 0.9846 with tau\_min = 0.734 across four diverse calibration distributions.
\item Threshold invariance: all six cross-epsilon Kendall's tau values exceed 0.95 across epsilon in {0.001, 0.01, 0.05, 0.1} (minimum: tau = 0.9597).
\item An empirical foundation for treating the sparsity fingerprint as a zero-cost structural prior for LoRA rank allocation, along with a clear statement of what remains unvalidated.
\end{enumerate}

